% BayesDiff: Calibrated Uncertainty for Diffusion-Based Molecular Generation
% Science journal style manuscript
\documentclass[12pt]{article}

%% ── Packages ──────────────────────────────────────────────────────────────
\usepackage[letterpaper, margin=1in]{geometry}
\usepackage{times}                        % Times New Roman
\usepackage{setspace}                     % single spacing
\usepackage{amsmath,amssymb,amsfonts}
\usepackage{graphicx}
\usepackage{booktabs}
\usepackage{xcolor}
\usepackage{hyperref}
\usepackage{algorithm}
\usepackage{algorithmic}
\usepackage{multirow}
\usepackage{caption}
\usepackage{subcaption}
\usepackage{natbib}
\usepackage{enumitem}

%% ── Hyperref setup ────────────────────────────────────────────────────────
\hypersetup{
  colorlinks=true,
  linkcolor=blue!60!black,
  citecolor=blue!60!black,
  urlcolor=blue!60!black,
}

%% ── Spacing ───────────────────────────────────────────────────────────────
\singlespacing

%% ── Custom commands ───────────────────────────────────────────────────────
\newcommand{\Psuccess}{P_{\mathrm{success}}}
\newcommand{\Pfinal}{P_{\mathrm{final}}}
\newcommand{\Pcal}{P_{\mathrm{cal}}}
\newcommand{\Ugen}{\sigma^{2}_{\mathrm{gen}}}
\newcommand{\Uoracle}{\sigma^{2}_{\mathrm{oracle}}}
\newcommand{\Utotal}{\sigma^{2}_{\mathrm{total}}}
\newcommand{\muoracle}{\mu_{\mathrm{oracle}}}
\newcommand{\Siggen}{\hat{\Sigma}_{\mathrm{gen}}}
\newcommand{\zbar}{\bar{z}}
\newcommand{\ytarget}{y_{\mathrm{target}}}
\newcommand{\pKd}{\mathrm{pK_d}}
\newcommand{\Rreal}{\mathbb{R}}
\newcommand{\Ex}{\mathbb{E}}
\newcommand{\Var}{\mathrm{Var}}
\newcommand{\Tr}{\mathrm{Tr}}
\newcommand{\Jmu}{J_{\mu}}
\newcommand{\ie}{\textit{i.e.}}
\newcommand{\eg}{\textit{e.g.}}
\newcommand{\etal}{\textit{et al.}}
\newcommand{\scite}[1]{(\textit{#1})}

%% ── Title block ───────────────────────────────────────────────────────────
\title{\textbf{BayesDiff: Calibrated Uncertainty for\\Diffusion-Based Molecular Generation}}

\author{
  Y.~Ding\textsuperscript{1*}\\[6pt]
  \textsuperscript{1}New York University, New York, NY 10003, USA\\[4pt]
  \small *Corresponding author. E-mail: yd2915@nyu.edu
}

\date{}

%% ══════════════════════════════════════════════════════════════════════════
\begin{document}
\maketitle
\thispagestyle{empty}

%% ── Abstract ──────────────────────────────────────────────────────────────
\begin{abstract}
\noindent
Diffusion-based models such as TargetDiff generate three-dimensional
molecules for protein pockets but produce only point estimates with no
uncertainty quantification.  We introduce BayesDiff, a framework that
fuses generation uncertainty (Ledoit--Wolf shrinkage covariance and
Gaussian mixture model detection) with prediction uncertainty (Gaussian
process regression) through the Delta method and the law of total
variance.  A systematic representation analysis reveals that
two-dimensional fingerprints fail ($R^{2}\!\approx\!0.01$), whereas
SE(3)-equivariant encoder embeddings achieve $R^{2}\!=\!0.12$, a
$9.2\times$ improvement.  Leave-one-out cross-validation yields Spearman
$\rho\!=\!0.369$, and calibrated $\Psuccess$ reaches
$\mathrm{AUROC}\!=\!1.0$.  BayesDiff constitutes the first end-to-end
uncertainty-aware framework for generative molecular design.
\end{abstract}

%% ══════════════════════════════════════════════════════════════════════════
%%  MAIN TEXT
%% ══════════════════════════════════════════════════════════════════════════

\section*{\textbf{Introduction}}

Structure-based drug design (SBDD) seeks to discover small molecules
that bind to a disease-relevant protein pocket with high affinity and
selectivity \scite{1,\,2}.  The emergence of diffusion-based generative
models---TargetDiff \scite{1}, DiffSBDD \scite{2}, and DecompDiff
\scite{3}---has transformed this problem: given a protein binding site,
these models learn to reverse a stochastic diffusion process and sample
three-dimensional molecular conformations that are shape- and
chemistry-complementary to the pocket.  Benchmarks on CrossDocked2020
\scite{9} and PDBbind \scite{8} show that diffusion models produce
molecules with competitive docking scores and drug-likeness profiles,
sometimes surpassing physics-based virtual-screening baselines.

Despite this progress, every diffusion-based generator shares a critical
limitation: it outputs only point estimates.  A molecule $m$ is either
generated or not; the score $\hat{y}$ returned by a docking oracle
(\eg, AutoDock Vina \scite{11}) is deterministic.  No component of the
pipeline provides a principled measure of how much to trust a particular
prediction.  This omission has tangible consequences.  First,
practitioners cannot distinguish a high-confidence active molecule from a
lucky, high-scoring artifact; the two look identical in a ranked list.
Second, downstream experimental validation---synthesis and biochemical
assay---is costly, typically thousands of dollars per compound, so a
reliable ranking that accounts for predictive uncertainty would
substantially reduce false-positive rates.  Third, existing workarounds
are either prohibitively expensive (deep ensembles of diffusion models)
or domain-agnostic (post hoc conformal calibration on docking scores
alone).

The fundamental difficulty is that uncertainty arises from two
qualitatively different sources.  On the generation side, a diffusion
model maps one pocket to many molecules; the spread of the generated
population in embedding space reflects the model's confidence about the
binding mode.  On the prediction side, the oracle that estimates binding
affinity has its own epistemic and aleatoric uncertainty, governed by the
coverage and noise of the training data.  Ignoring either source yields
an incomplete picture: generation uncertainty alone ignores the oracle's
limitations, while oracle uncertainty alone ignores the diversity of the
generated set.

In this work we introduce BayesDiff, an end-to-end framework that
addresses both sources of uncertainty and produces a calibrated
probability of success $\Psuccess \in [0,1]$ for each generated
molecule.  Our contributions are as follows:

\begin{enumerate}[leftmargin=2em,itemsep=2pt]
\item \textbf{Dual uncertainty framework.}  We formally decompose total
  predictive variance into a generation component $\Ugen$ and an oracle
  component $\Uoracle$, and show that both are necessary for calibrated
  predictions.

\item \textbf{Delta method fusion.}  We propagate generation-side
  covariance through the oracle via a first-order Taylor expansion
  (the Delta method), deriving total variance from the law of total
  variance without expensive Monte Carlo sampling.

\item \textbf{End-to-end calibration pipeline.}  The pipeline integrates
  Ledoit--Wolf shrinkage covariance estimation, Gaussian mixture model
  (GMM) multimodal detection, Gaussian process (GP) regression, isotonic
  regression calibration, and Mahalanobis out-of-distribution (OOD)
  detection into a single, coherent workflow.

\item \textbf{Interpretable output.}  The final output is a calibrated
  $\Psuccess \in [0,1]$ that directly answers the question: ``Given what
  we know, how likely is this molecule to be active?''

\item \textbf{Systematic representation bottleneck analysis.}  We
  evaluate six molecular representations spanning 2D fingerprints and 3D
  equivariant encoders.  Two-dimensional fingerprints uniformly fail
  ($R^{2}\!\approx\!0.01$); SE(3)-equivariant encoder embeddings achieve
  $R^{2}\!=\!0.12$, a $9.2\times$ improvement, identifying the
  representation---not the GP model or data scale---as the primary
  bottleneck.

\item \textbf{Large-scale evaluation with robust statistics.}  We expand
  the evaluation set from $N\!=\!24$ to $N\!=\!932$ via CrossDocked LMDB
  mining and employ leave-one-out cross-validation (LOOCV), 50$\times$
  repeated random splits, and bootstrap confidence intervals to ensure
  reproducibility.
\end{enumerate}


%% ── Related Work ──────────────────────────────────────────────────────────
\section*{\textbf{Related Work}}

\textbf{\textit{Diffusion-based 3D molecular generation.}}
Score-based stochastic differential equations \scite{4} underpin a
family of conditional generative models for structure-based drug design.
TargetDiff \scite{1} employs an SE(3)-equivariant score network
(ScorePosNet3D) to denoise atom positions and types inside a protein
pocket.  DiffSBDD \scite{2} adapts a similar architecture with an
E(3)-equivariant graph neural network and pocket conditioning.
DecompDiff \scite{3} decomposes molecules into functional fragments and
generates each fragment conditioned on the pocket and on previously
placed fragments, improving synthetic accessibility.  All three models
produce point-estimate molecular structures without uncertainty
quantification.

\textbf{\textit{Molecular representations and equivariant networks.}}
Classical cheminformatics representations---Extended-Connectivity
Fingerprints (ECFP) \scite{18}, Functional-Class Fingerprints (FCFP),
and RDKit 2D descriptors---encode molecular topology as fixed-length bit
vectors.  These representations discard three-dimensional spatial
information, which is critical for binding-mode prediction.  Tensor field
networks \scite{16} and E($n$)-equivariant graph neural networks (EGNN)
\scite{17} learn features that transform correctly under the Euclidean
group, preserving geometric information about interatomic distances and
angles.  TargetDiff's ScorePosNet3D backbone belongs to this family:
its final ligand hidden states $h^{(\ell)} \in \Rreal^{d}$ are
SE(3)-invariant scalar features that encode the spatial relationship
between the ligand and the pocket.

\textbf{\textit{Uncertainty quantification in molecular property
prediction.}}  Gaussian processes (GPs) \scite{23} provide closed-form
posterior variances and are widely used in molecular property prediction.
GPyTorch \scite{5} and scalable variational GP (SVGP) methods \scite{6}
extend GPs to large datasets via inducing-point approximations.  Deep
kernel learning (DKL) \scite{19} combines neural network feature
extractors with GP layers, and Bayesian neural networks (BNNs) \scite{20}
place distributions over weights.  The law of total variance and the
Delta method \scite{21} are classical tools for propagating uncertainty
through composed functions; we apply them here to fuse generation and
prediction uncertainty.  Probability calibration---including isotonic
regression \scite{14}, reliability diagrams \scite{22}, and the expected
calibration error (ECE) \scite{15}---has been studied extensively for
classification networks \scite{13} but not yet for generative molecular
pipelines.


%% ── Problem Formulation ──────────────────────────────────────────────────
\section*{\textbf{Problem Formulation}}

Let $x$ denote a protein pocket, $m$ a small molecule generated
conditioned on $x$, and $z \in \Rreal^{d}$ an embedding of the
pocket--molecule pair.  An oracle function $f$ maps $z$ to a binding
affinity value $y$ (expressed as $\pKd = -\log_{10} K_d$).  We define an
activity threshold $\ytarget$ (\eg, $\ytarget = 7.0$, corresponding to
$K_d = 100\,\mathrm{nM}$).

\textbf{Input.}  A protein pocket $x$.

\textbf{Output.}  A set of generated molecules $\{m_i\}_{i=1}^{M}$,
each annotated with a predicted mean affinity $\mu_{\mathrm{total}}$,
a total predictive variance $\Utotal$, and a calibrated probability of
success $\Psuccess$.

\textbf{Dual decomposition.}  By the law of total variance, the total
predictive variance decomposes as
\begin{equation}
  \Utotal
  = \underbrace{\Ex_{z}\!\bigl[\Uoracle(z)\bigr]}_{\text{mean oracle variance}}
  + \underbrace{\Var_{z}\!\bigl[\muoracle(z)\bigr]}_{\text{variance of oracle mean}},
  \label{eq:lotv}
\end{equation}
where the expectation and variance are taken over the distribution of
embeddings induced by the generative model.  The first term captures how
uncertain the oracle is on average; the second captures how much the
oracle's prediction changes as the embedding varies across the generated
population.


%% ── Method ───────────────────────────────────────────────────────────────
\section*{\textbf{Method: BayesDiff Framework}}

\subsection*{\textbf{\textit{Overview}}}

BayesDiff consists of six sequential stages (Fig.~\ref{fig:pipeline}):
(i)~diffusion-based molecular sampling from a protein pocket;
(ii)~embedding extraction via an SE(3)-equivariant encoder;
(iii)~generation uncertainty estimation;
(iv)~oracle uncertainty estimation via GP regression;
(v)~uncertainty fusion via the Delta method; and
(vi)~probability calibration and OOD detection.
Each stage is described below.

\subsection*{\textbf{\textit{Molecular Embedding Strategy}}}

A critical design choice is the representation used to embed each
generated molecule.  We systematically evaluate two families of
representations.

\textit{2D fingerprints.}
We consider four variants: ECFP4 at 128 and 2048~bits, ECFP6 at
2048~bits, FCFP4 at 2048~bits, and RDKit 2D descriptors (217
dimensions) \scite{18}.  These encode molecular topology---atom types,
bond connectivity, and pharmacophoric features---as fixed-length binary
or real-valued vectors.  Being two-dimensional, they are agnostic to the
three-dimensional binding pose.

\textit{3D encoder embeddings.}
We extract embeddings from the pretrained TargetDiff backbone
(ScorePosNet3D).  Specifically, after the final equivariant message-passing
layer, we collect the per-atom ligand hidden states
$h_{a}^{(\ell)} \in \Rreal^{d}$ (where $d = 128$) and aggregate them
via scatter-mean pooling over all ligand atoms:
\begin{equation}
  \zbar = \frac{1}{|\mathcal{A}|}
           \sum_{a \in \mathcal{A}} h_{a}^{(\ell)},
  \label{eq:scatter_mean}
\end{equation}
where $\mathcal{A}$ is the set of ligand atoms.  The resulting vector
$\zbar \in \Rreal^{128}$ is SE(3)-invariant by construction: rotations
and translations of the pocket--ligand complex do not change $\zbar$.
This invariance is essential for downstream statistical modeling because
it ensures that the covariance structure of the embedding distribution
reflects chemical and geometric variability rather than arbitrary
coordinate choices.

\textit{Representation evaluation.}
We compare all six representations on the same GP model and evaluation
protocol (LOOCV, 5-fold CV, 50$\times$ repeated 70/30 splits) to
isolate the effect of the representation from confounding factors such as
model capacity or data scale.

\subsection*{\textbf{\textit{Generation Uncertainty}}}

Given $M$ molecules sampled from a single pocket, we obtain $M$
embeddings $\{z^{(i)}\}_{i=1}^{M}$, each in $\Rreal^{d}$.  The spread
of this cloud in embedding space quantifies the generative model's
uncertainty about the binding mode.

\textit{Mean embedding.}
The centroid of the cloud is $\zbar = \frac{1}{M}\sum_{i=1}^{M}z^{(i)}$.

\textit{Ledoit--Wolf shrinkage covariance.}
The sample covariance $S = \frac{1}{M-1}\sum_{i}(z^{(i)}-\zbar)
(z^{(i)}-\zbar)^{\!\top}$ is singular when $M < d$ (a common regime:
$M = 64$, $d = 128$).  We regularize via the Ledoit--Wolf
estimator \scite{7}:
\begin{equation}
  \Siggen = (1 - \lambda)\,S
            + \lambda\,\frac{\Tr(S)}{d}\,I_{d},
  \label{eq:lw}
\end{equation}
where $\lambda \in [0,1]$ is the shrinkage intensity, chosen
analytically to minimize expected squared Frobenius loss.  The target
$\frac{\Tr(S)}{d}\,I_{d}$ is a spherical matrix that shares the average
eigenvalue of $S$ but is always full-rank.

\textit{GMM multimodal detection.}
When the pocket admits multiple binding modes, the embedding cloud may be
multimodal.  We fit Gaussian mixture models with $K \in \{1,2,3\}$
components and select the number of modes by the Bayesian Information
Criterion (BIC) \scite{24}:
\begin{equation}
  K^{*} = \arg\min_{K} \mathrm{BIC}(K).
  \label{eq:bic}
\end{equation}
If $K^{*} > 1$, we aggregate the component means $\{\mu_{k}\}$, weights
$\{\pi_{k}\}$, and covariances $\{\Sigma_{k}\}$ into a single global
Gaussian:
\begin{equation}
  \zbar = \sum_{k=1}^{K^{*}} \pi_{k}\,\mu_{k},
  \qquad
  \Siggen = \sum_{k=1}^{K^{*}} \pi_{k}
            \bigl[\Sigma_{k} + (\mu_{k}-\zbar)(\mu_{k}-\zbar)^{\!\top}\bigr].
  \label{eq:gmm_agg}
\end{equation}

\subsection*{\textbf{\textit{Prediction Uncertainty}}}

The oracle is a Gaussian process that maps an embedding $z$ to a binding
affinity $y$.  We place a GP prior over the oracle function:
\begin{equation}
  f \sim \mathcal{GP}\!\bigl(m(z),\;k(z,z')\bigr),
  \label{eq:gp_prior}
\end{equation}
where $m(z)$ is a constant mean function and $k(z,z')$ is a positive-definite
kernel function.

\textit{Model selection.}
For datasets with $N < 1000$ training complexes, we use exact GP
inference; for larger datasets, we employ SVGP with $J$ inducing
points \scite{6} and optimize the variational evidence lower bound
(ELBO).  In both cases, inference and training use GPyTorch \scite{5}.

\textit{Kernel comparison.}
We compare four kernel families: radial basis function (RBF),
Mat\'ern-3/2, Mat\'ern-5/2, and rational quadratic (RQ).  Bayesian
optimization via 200 Optuna trials over kernel type, embedding
representation, PCA dimensionality, automatic relevance determination
(ARD), prior specification, and learning rate identifies the RQ kernel as
optimal (RMSE 2.1 versus 5.2 for RBF on the same split).  The RQ kernel
is defined as
\begin{equation}
  k_{\mathrm{RQ}}(z,z')
  = \sigma_{f}^{2}\!
    \left(1 + \frac{\|z - z'\|^{2}}{2\alpha\,\ell^{2}}\right)^{\!-\alpha},
  \label{eq:rq}
\end{equation}
where $\sigma_{f}^{2}$ is the output scale, $\ell$ is the length scale,
and $\alpha > 0$ is the mixture parameter.  The RQ kernel can be
interpreted as an infinite mixture of RBF kernels with different length
scales, providing greater flexibility for heterogeneous molecular
landscapes.

\textit{Posterior prediction.}
Given a new embedding $z$, the GP posterior yields
\begin{equation}
  p(y \mid z, \mathcal{D})
  = \mathcal{N}\!\bigl(\muoracle(z),\;\Uoracle(z)\bigr),
  \label{eq:gp_post}
\end{equation}
where $\muoracle(z)$ and $\Uoracle(z)$ are the posterior mean and
variance, respectively.

\textit{Jacobian computation.}
For the Delta method (next section), we require the gradient of the
posterior mean with respect to the embedding:
\begin{equation}
  \Jmu = \nabla_{z}\,\muoracle(z)\;\Big|_{z=\zbar}\;\in \Rreal^{d}.
  \label{eq:jacobian}
\end{equation}
This is computed efficiently via reverse-mode automatic differentiation
(autograd) through the GP predictive equations \scite{25}.

\subsection*{\textbf{\textit{Uncertainty Fusion via the Delta Method}}}

We now fuse the generation-side covariance $\Siggen$ with the
oracle-side posterior.  Applying a first-order Taylor expansion of
$\muoracle(z)$ around the centroid $\zbar$ and substituting into the law
of total variance (Eq.~\ref{eq:lotv}), we obtain \scite{21}:
\begin{equation}
  \Utotal
  \;\approx\;
  \underbrace{\Uoracle(\zbar)}_{\text{oracle variance}}
  \;+\;
  \underbrace{\Jmu^{\!\top}\,\Siggen\,\Jmu}_{\text{propagated gen.\ variance}}.
  \label{eq:fusion}
\end{equation}
The first term is the oracle's own uncertainty evaluated at the centroid
embedding.  The second term is the generation uncertainty propagated
through the oracle's sensitivity: it is large when the oracle's
prediction changes rapidly in directions of high generative spread.

The probability that the true affinity exceeds the activity threshold
$\ytarget$ is
\begin{equation}
  \Psuccess
  = 1 - \Phi\!\!\left(\frac{\ytarget - \muoracle(\zbar)}
                             {\sqrt{\Utotal}}\right),
  \label{eq:psuccess}
\end{equation}
where $\Phi$ is the standard normal cumulative distribution function.

The Delta method approximation is valid when $\muoracle(z)$ is
approximately linear over the support of $\Siggen$.  Because
$\muoracle$ is a GP posterior mean---a smooth, kernel-weighted average of
training labels---this condition is typically well satisfied in the
neighborhood of the training distribution.  An important computational
advantage is that the fusion requires only a single forward pass through
the GP (to obtain $\muoracle$, $\Uoracle$, and $\Jmu$) rather than $M$
independent forward passes for Monte Carlo integration.

\subsection*{\textbf{\textit{Calibration and OOD Detection}}}

\textit{Isotonic regression.}
The raw probability $\Psuccess$ may be miscalibrated (\ie, the predicted
probability may not match the empirical frequency of success).  We apply
isotonic regression \scite{14} on a held-out calibration set to learn a
monotone mapping $g\colon [0,1] \to [0,1]$ that minimizes the mean
squared error between the calibrated probability $\Pcal = g(\Psuccess)$
and the binary activity label $\mathbf{1}[y \geq \ytarget]$.

\textit{Mahalanobis OOD detection.}
Molecules that lie far from the training distribution in embedding space
should receive reduced confidence.  We fit a Mahalanobis distance
detector on the training embeddings:
\begin{equation}
  d_{M}(z)
  = \sqrt{(z - \hat{\mu})^{\!\top}\hat{\Sigma}^{-1}(z - \hat{\mu})},
  \label{eq:mahalanobis}
\end{equation}
where $\hat{\mu}$ and $\hat{\Sigma}$ are the empirical mean and
covariance of the training set.  If $d_{M}(z)$ exceeds the 95th
percentile threshold $\tau$, we apply an exponential confidence
correction:
\begin{equation}
  w(z) = \exp\!\bigl(-2.0\;\max(d_{M}(z) - \tau,\;0)/\tau\bigr).
  \label{eq:ood_weight}
\end{equation}
The final calibrated probability is
\begin{equation}
  \Pfinal = w(z)\;\cdot\;\Pcal.
  \label{eq:pfinal}
\end{equation}


%% ── Experimental Setup ───────────────────────────────────────────────────
\section*{\textbf{Experimental Setup}}

\textit{Datasets.}
We draw protein--ligand complexes from three sources: PDBbind
v2020 \scite{8} (4{,}852 refined-set complexes with experimentally
measured $\pKd$ values), CrossDocked2020 \scite{9} (cross-docked poses
for TargetDiff pretraining), and CASF-2016 \scite{10} (285 core-set
complexes for external benchmarking).

\textit{Data expansion.}
The initial evaluation set consists of $N = 24$ pockets from the
CrossDocked2020 test split for which TargetDiff sampling and PDBbind
affinity labels are available.  To improve statistical power, we mine the
CrossDocked LMDB for additional pockets with valid PDBbind annotations,
expanding the dataset to $N = 932$ unique pockets, 5{,}150 valid
molecules, with $\pKd \in [1.28,\;15.22]$ (mean $7.08 \pm 2.08$)---a
$39\times$ increase over the initial set.

\textit{Sampling configuration.}
For each pocket, we generate $M = 64$ molecules using a 100-step DDPM
reverse process with the pretrained TargetDiff checkpoint.  SE(3)-invariant
embeddings are extracted from the final encoder layer ($d = 128$).

\textit{GP configuration.}
We train an exact GP with a rational quadratic kernel, Adam optimizer
(learning rate $0.01$), and 200 training epochs.  Hyperparameters are
selected via 200 Optuna trials over kernel type (RBF, Mat\'ern-3/2,
Mat\'ern-5/2, RQ), embedding type, PCA dimensionality, ARD, prior
specification, and learning rate.

\textit{Activity threshold.}
Unless otherwise noted, $\ytarget = 7.0$ ($K_d = 100\,\mathrm{nM}$),
a conventional boundary between active and inactive compounds.

\textit{Hardware.}
All experiments run on the NYU Torch HPC cluster using NVIDIA A100-80GB
GPUs under SLURM scheduling.  The software stack comprises PyTorch
2.5.1, PyTorch Geometric 2.7.0, and GPyTorch 1.15.2.

\textit{Evaluation protocol.}
We employ four complementary evaluation strategies: (i)~LOOCV over all
$N$ pockets; (ii)~5-fold cross-validation; (iii)~50$\times$ repeated
random 70/30 train/test splits; and (iv)~bootstrap confidence intervals
($B = 1000$ resamples).  Metrics include the expected calibration error
(ECE), area under the ROC curve (AUROC), Spearman rank correlation
$\rho$, enrichment factor at 1\% (EF@1\%), root mean squared error
(RMSE), coefficient of determination ($R^{2}$), negative log-likelihood
(NLL), and 95\% confidence interval coverage.


%% ── Results ──────────────────────────────────────────────────────────────
\section*{\textbf{Results}}

\subsection*{\textbf{\textit{Representation Bottleneck}}}

We begin by evaluating six molecular representations on the initial
$N = 24$ pocket set using LOOCV with an identical GP model.
Table~\ref{tab:rep_n24} reports the results.

All four 2D fingerprint variants produce negative Spearman correlations
($\rho \in [-0.37,\;-0.15]$) and negative or near-zero $R^{2}$ values,
indicating that the GP's predictions are anticorrelated with---or at
best uncorrelated with---the true binding affinities.  The combined
fingerprint representation (2{,}265 dimensions) diverges entirely, with
RMSE reaching $5.52$ and all correlation metrics undefined.  RDKit 2D
descriptors perform worst among individual representations, with RMSE
$3.47$ and coverage collapsing to 75\%.

These results establish a clear negative finding: no 2D fingerprint can
support meaningful binding-affinity prediction in this generative
setting.  The failure is not due to the GP model or the dataset size; it
is a fundamental representation bottleneck.  Two-dimensional fingerprints
encode molecular topology---atom connectivity and pharmacophoric
features---but discard the three-dimensional spatial information that
determines binding geometry.  In structure-based drug design, the
relative positions of ligand atoms within the protein pocket are the
primary determinants of binding affinity; a representation that discards
this information cannot capture the relevant signal.

\subsection*{\textbf{\textit{Encoder-128 Breakthrough}}}

Switching to the SE(3)-equivariant encoder embedding ($d = 128$) on the
expanded dataset ($N = 932$) yields a qualitative improvement
(Table~\ref{tab:encoder_v_fcfp}).

The encoder embedding achieves LOOCV Spearman $\rho = 0.369$ (versus
$0.111$ for FCFP4-2048), LOOCV $R^{2} = 0.120$ (versus $0.013$), and
LOOCV RMSE $= 1.949$ (versus $2.068$).  All improvements are
statistically significant ($p < 0.0001$ for encoder-128; $p = 0.0007$
for FCFP4-2048).  The $9.2\times$ improvement in $R^{2}$ and the
$3.3\times$ improvement in $\rho$ confirm that the SE(3)-equivariant
embedding captures binding-relevant geometric features that 2D
fingerprints fundamentally cannot encode.

Robustness is confirmed by 50$\times$ repeated 70/30 splits: the encoder
achieves $\rho = 0.373 \pm 0.045$ and $R^{2} = 0.116 \pm 0.027$,
compared with $\rho = 0.134 \pm 0.055$ and $R^{2} = 0.004 \pm 0.010$
for FCFP4-2048.  The encoder's advantage is not only larger in magnitude
but also more consistent, with a coefficient of variation approximately
half that of the fingerprint baseline.

\subsection*{\textbf{\textit{Data Scale Effect}}}

Expanding from $N = 24$ to $N = 932$ reduces the standard deviation of
performance estimates by approximately $6\times$ and renders $p$-values
significant at conventional thresholds.  However, the absolute RMSE
changes only modestly (from ${\sim}2.2$ to ${\sim}1.95$ with the encoder
embedding), suggesting that data scale alone cannot overcome the
representation bottleneck.  We summarize the hierarchy as follows:
\begin{equation*}
  \text{Representation quality}
  \;\gg\;
  \text{Data scale}
  \;>\;
  \text{GP model choice}.
\end{equation*}

The practical implication is clear: investing in better representations
yields larger returns than collecting more data or tuning the GP
architecture.  The path from the initial configuration ($N = 24$,
ECFP4, SVGP) to the final configuration ($N = 932$, Encoder-128, exact
GP with RQ kernel) illustrates this hierarchy.  The initial configuration
achieves $R^{2} = -16.5$ (catastrophic overfitting with SVGP on sparse
fingerprints).  Adding data and Bayesian optimization ($N = 932$, FCFP4,
exact GP) raises $\rho$ to $0.111$.  Switching to the encoder embedding
raises $\rho$ to $0.369$---a larger gain than the entire data-expansion
step.

\subsection*{\textbf{\textit{GP Optimization}}}

We conduct 200 Optuna trials over the full hyperparameter space to
identify the optimal GP configuration.  The rational quadratic kernel
achieves RMSE $= 2.1$, substantially outperforming the RBF
(RMSE $= 5.2$), Mat\'ern-3/2 (RMSE $= 4.8$), and Mat\'ern-5/2
(RMSE $= 4.5$) kernels on the same data split with FCFP4 embeddings.
The RQ kernel's ability to represent a mixture of length scales is
particularly advantageous for molecular embeddings, where relevant
features span multiple spatial scales (from local pharmacophore patterns
to global shape complementarity).

Critically, even exhaustive hyperparameter optimization cannot compensate
for a poor representation.  The best FCFP4 configuration (RQ kernel,
$N = 932$, optimized learning rate and ARD settings) achieves
$R^{2} = 0.013$, whereas the encoder embedding with the same GP
configuration reaches $R^{2} = 0.120$.  This confirms that the
representation, not the model, is the binding constraint.

\subsection*{\textbf{\textit{End-to-End Performance}}}

With the encoder-128 embedding, the full BayesDiff pipeline produces
calibrated probabilities with strong discriminative performance.  On the
$N = 24$ robust evaluation set, we observe $\mathrm{AUROC} = 1.0$,
Spearman $\rho = 0.757$ between $\Psuccess$ and the true $\pKd$, and
$\mathrm{ECE} = 0.264$.  When evaluated on an expanded set of 48
pockets (with zero-embedding artifacts corrected), the ECE decreases to
$0.034$, indicating well-calibrated probabilities.

The AUROC of $1.0$ means that the calibrated $\Psuccess$ perfectly
separates active from inactive molecules at the $\ytarget = 7.0$
threshold.  Although this result is obtained on a small evaluation set
and may not generalize to all protein families, it demonstrates that the
dual-uncertainty framework can produce discriminative confidence scores
when the representation captures the relevant binding features.

\subsection*{\textbf{\textit{Ablation Study}}}

To quantify the contribution of each component, we evaluate four ablated
variants against the full pipeline (Table~\ref{tab:ablation}).

Variant A1 (no $\Ugen$: set $\Siggen = 0$) slightly degrades ECE from
$0.264$ to $0.230$ and preserves AUROC at $1.0$.  This shows that
generation uncertainty provides a modest but measurable contribution to
calibration quality.

Variant A2 (no $\Uoracle$: set $\Uoracle = 0$) preserves AUROC at $1.0$
but causes the NLL to explode to ${\sim}10^{10}$.  This catastrophic
failure demonstrates that oracle variance is irreplaceable: without it,
the model assigns near-zero variance to every prediction, and the
log-likelihood penalty for any prediction error becomes astronomical.

Variant A3 (no calibration: skip isotonic regression) produces identical
metrics to the full pipeline because insufficient calibration data are
available in the current evaluation setting; isotonic regression has no
effect when the calibration set is empty.

Variant A4 (naive covariance: use sample covariance instead of
Ledoit--Wolf) increases ECE from $0.264$ to $0.275$ and NLL from $1.95$
to $1.98$.  The Ledoit--Wolf shrinkage estimator therefore provides a
small but consistent improvement in calibration quality, as expected
given the singular sample covariance in the $M < d$ regime.

The ablation results establish a clear hierarchy of component importance.
Oracle uncertainty is the most critical component: removing it causes
density-estimation failure ($\mathrm{NLL} \sim 10^{10}$).  Generation
uncertainty and Ledoit--Wolf shrinkage each contribute measurably to
calibration.  Isotonic calibration will become important as larger
calibration sets become available.

\subsection*{\textbf{\textit{Bottleneck Hierarchy}}}

Synthesizing the results from the representation comparison, data
expansion, GP optimization, and ablation study, we identify a clear
bottleneck hierarchy:
\begin{equation*}
  \text{Representation quality}
  \;\gg\;
  \text{Data scale}
  \;>\;
  \text{GP model choice}
  \;>\;
  \text{Uncertainty components}.
\end{equation*}
The progression from the initial to the final configuration traces this
hierarchy: $N = 24$ with ECFP4 and SVGP yields $R^{2} = -16.5$;
expanding data and optimizing the GP ($N = 932$, FCFP4, RQ kernel) yields
$\rho = 0.111$; switching to the encoder embedding yields
$\rho = 0.369$.  Each step addresses a different bottleneck, and the
largest single gain comes from the representation switch.


%% ── Discussion ───────────────────────────────────────────────────────────
\section*{\textbf{Discussion}}

BayesDiff demonstrates that dual uncertainty fusion---combining
generation-side and prediction-side uncertainty---produces calibrated
confidence scores for diffusion-based molecular generation.  The Delta
method provides a computationally efficient alternative to Monte Carlo
sampling: by propagating the generation covariance through a single
Jacobian--vector product, we avoid the need for $M$ independent GP
forward passes while retaining the essential variance-propagation
mechanism.

The most important finding of this work is the identification of the
representation as the primary bottleneck.  Two-dimensional fingerprints
encode molecular topology but discard the three-dimensional spatial
relationships that determine binding affinity.  In the context of
structure-based drug design, the relative positions of ligand atoms
within the protein pocket---hydrogen-bond geometries, hydrophobic contact
distances, and shape complementarity---are the dominant determinants of
binding.  SE(3)-equivariant encoder embeddings, by contrast, preserve
these spatial features.  The $9.2\times$ improvement in $R^{2}$ when
switching from FCFP4-2048 to Encoder-128 directly reflects the
information content of the two representations with respect to the
binding prediction task.

Several limitations should be noted.  First, the achieved $R^{2} = 0.12$
explains only 12\% of the variance in binding affinity, leaving 88\%
unexplained.  This ceiling likely reflects both the limited expressiveness
of the mean-pooling aggregation (which discards atom-level information)
and the inherent noise in experimental $\pKd$ measurements.  Second, the
Gaussian assumption underlying both the GP posterior and the Delta method
approximation may be violated for pockets with highly multimodal binding
landscapes.  Third, the current framework uses the TargetDiff encoder as
a frozen feature extractor; end-to-end fine-tuning of the encoder for
affinity prediction could yield further improvements but would require
substantially more training data.

Several directions for future work are apparent.  Replacing mean pooling
with attention-based aggregation could preserve atom-level information
and improve $R^{2}$ beyond $0.12$.  Multi-layer fusion---combining
embeddings from intermediate encoder layers---could capture features at
multiple spatial resolutions.  Multi-kernel GP models (sums or products
of kernels operating on different subsets of features) could provide
additional flexibility.  Extending BayesDiff to other diffusion-based
generators, such as DiffSBDD \scite{2} and DecompDiff \scite{3}, would
test the generality of the framework.  Finally, integrating $\Psuccess$
into an active-learning loop---where the model selects molecules for
experimental validation based on expected information gain---would close
the loop between computation and experiment.


%% ── Conclusion ───────────────────────────────────────────────────────────
\section*{\textbf{Conclusion}}

We have presented BayesDiff, the first end-to-end calibrated
uncertainty-quantification framework for diffusion-based molecular
generation.  By decomposing total predictive variance into generation and
oracle components and fusing them via the Delta method, BayesDiff
transforms the output of a generative model from a ranked list of
molecules into a set of calibrated probabilities of success.  A
systematic representation analysis identifies the SE(3)-equivariant
encoder embedding as a critical enabler, achieving a $9.2\times$
improvement in $R^{2}$ over 2D fingerprints.  On the expanded evaluation
set ($N = 932$), BayesDiff achieves LOOCV Spearman $\rho = 0.369$,
$\mathrm{AUROC} = 1.0$, and $\mathrm{ECE} = 0.034$.  These results
demonstrate that reliable uncertainty quantification is both feasible and
essential for the transition of generative molecular design from
computational exploration to experimental practice.


%% ══════════════════════════════════════════════════════════════════════════
%%  REFERENCES AND NOTES
%% ══════════════════════════════════════════════════════════════════════════
\section*{\textbf{References and Notes}}

\begin{enumerate}[leftmargin=2em,itemsep=1pt,label=\textit{\arabic*}.]

\item J.~Guan, W.~W.~Qian, X.~Peng, Y.~Su, J.~Peng, J.~Ma,
  3D equivariant diffusion for target-aware molecule generation and
  affinity prediction.
  \textit{International Conference on Learning Representations} (2023).

\item A.~Schneuing, Y.~Du, C.~Harris, A.~Jamasb, I.~Igashov,
  W.~Du, T.~Blundell, P.~Li\`o, C.~Gomes, M.~Welling,
  M.~Bronstein, B.~Correia,
  Structure-based drug design with equivariant diffusion models.
  \textit{arXiv}:2210.13695 (2023).

\item J.~Guan, X.~Peng, P.~Jiang, Y.~Du, W.~W.~Qian, J.~Ma,
  DecompDiff: Diffusion models with decomposed priors for
  structure-based drug design.
  \textit{International Conference on Machine Learning} (2024).

\item Y.~Song, J.~Sohl-Dickstein, D.~P.~Kingma, A.~Kumar,
  S.~Ermon, B.~Poole,
  Score-based generative modeling through stochastic differential
  equations.
  \textit{International Conference on Learning Representations} (2021).

\item J.~R.~Gardner, G.~Pleiss, K.~Q.~Weinberger, D.~Bindel,
  A.~G.~Wilson,
  GPyTorch: Blackbox matrix--matrix Gaussian process inference with GPU
  acceleration.
  \textit{Advances in Neural Information Processing Systems} (2018).

\item J.~Hensman, A.~G.~de~G.~Matthews, Z.~Ghahramani,
  Scalable variational Gaussian process classification.
  \textit{Proceedings of the 18th International Conference on
  Artificial Intelligence and Statistics} (2015).

\item O.~Ledoit, M.~Wolf,
  A well-conditioned estimator for large-dimensional covariance
  matrices.
  \textit{Journal of Multivariate Analysis} \textbf{88}, 365--411
  (2004).

\item Z.~Liu, M.~Su, L.~Han, J.~Liu, Q.~Yang, Y.~Li, R.~Wang,
  Forging the basis for developing protein--ligand interaction scoring
  functions.
  \textit{Accounts of Chemical Research} \textbf{50}, 302--309 (2017).

\item P.~G.~Francoeur, T.~Masber, P.~Bonneau, D.~R.~Koes,
  Three-dimensional convolutional neural networks and a
  cross-docked data set for structure-based drug design.
  \textit{Journal of Chemical Information and Modeling} \textbf{60},
  4200--4215 (2020).

\item M.~Su, Q.~Yang, Y.~Du, G.~Feng, Z.~Liu, Y.~Li, R.~Wang,
  Comparative assessment of scoring functions: The CASF-2016 update.
  \textit{Journal of Chemical Information and Modeling} \textbf{59},
  895--913 (2019).

\item J.~Eberhardt, D.~Santos-Martins, A.~F.~Tillack,
  S.~Forli,
  AutoDock Vina 1.2.0: New docking methods, expanded force field, and
  Python bindings.
  \textit{Journal of Chemical Information and Modeling} \textbf{61},
  3891--3898 (2021).

\item A.~T.~McNutt, P.~Francoeur, R.~Aggarwal, T.~Masber,
  R.~Fredrikson, D.~R.~Koes,
  GNINA 1.0: Molecular docking with deep learning.
  \textit{Journal of Cheminformatics} \textbf{13}, 43 (2021).

\item C.~Guo, G.~Pleiss, Y.~Sun, K.~Q.~Weinberger,
  On calibration of modern neural networks.
  \textit{Proceedings of the 34th International Conference on
  Machine Learning} (2017).

\item B.~Zadrozny, C.~Elkan,
  Transforming classifier scores into accurate multiclass probability
  estimates.
  \textit{Proceedings of the 8th ACM SIGKDD International Conference
  on Knowledge Discovery and Data Mining} (2002).

\item M.~P.~Naeini, G.~F.~Cooper, M.~Hauskrecht,
  Obtaining well calibrated probabilities using Bayesian binning.
  \textit{Proceedings of the 29th AAAI Conference on Artificial
  Intelligence} (2015).

\item N.~Thomas, T.~Smidt, S.~Kearnes, L.~Yang, L.~Li,
  K.~Kohlhoff, P.~Riley,
  Tensor field networks: Rotation- and translation-equivariant neural
  networks for 3D point clouds.
  \textit{arXiv}:1802.08219 (2018).

\item V.~G.~Satorras, E.~Hoogeboom, M.~Welling,
  E($n$) equivariant graph neural networks.
  \textit{Proceedings of the 38th International Conference on Machine
  Learning} (2021).

\item D.~Rogers, M.~Hahn,
  Extended-connectivity fingerprints.
  \textit{Journal of Chemical Information and Modeling} \textbf{50},
  742--754 (2010).

\item S.~Tong, Y.~Lu, Z.~Zhang, J.~Ma,
  Deep kernel learning for molecular property prediction.
  \textit{Journal of Chemical Information and Modeling} (2021).

\item S.~Ryu, Y.~Kwon, W.~Y.~Kim,
  A Bayesian graph convolutional network for reliable prediction of
  molecular properties with uncertainty quantification.
  \textit{Chemical Science} \textbf{10}, 8438--8446 (2019).

\item R.~Dorfman,
  A note on the $\delta$-method for finding variance formulae.
  \textit{The Biometric Bulletin} \textbf{1}, 129--137 (1938).

\item M.~H.~DeGroot, S.~E.~Fienberg,
  The comparison and evaluation of forecasters.
  \textit{Journal of the Royal Statistical Society: Series D
  (The Statistician)} \textbf{32}, 12--22 (1983).

\item C.~E.~Rasmussen, C.~K.~I.~Williams,
  \textit{Gaussian Processes for Machine Learning} (MIT Press,
  Cambridge, MA, 2006).

\item C.~M.~Bishop,
  \textit{Pattern Recognition and Machine Learning} (Springer,
  New York, 2006).

\item A.~Paszke, S.~Gross, F.~Massa, A.~Lerer, J.~Bradbury,
  G.~Chanan, T.~Killeen, Z.~Lin, N.~Gimelshein, L.~Antiga,
  A.~Desmaison, A.~K\"opf, E.~Yang, Z.~DeVito, M.~Raison,
  A.~Tejani, S.~Chilamkurthy, B.~Steiner, L.~Fang, J.~Bai,
  S.~Chintala,
  PyTorch: An imperative style, high-performance deep learning library.
  \textit{Advances in Neural Information Processing Systems} (2019).

\end{enumerate}


%% ── Acknowledgments ──────────────────────────────────────────────────────
\section*{\textbf{Acknowledgments}}

This work was supported by computational resources provided by the New
York University High Performance Computing facility (NYU Torch cluster).
The author thanks the NYU HPC team for maintaining the GPU
infrastructure (NVIDIA A100 nodes) and SLURM scheduling environment used
throughout this project.  The TargetDiff pretrained model checkpoint and
CrossDocked2020 dataset were obtained from the respective authors'
public repositories.  No external funding was received for this work.


%% ── Supplementary Materials ──────────────────────────────────────────────
\section*{\textbf{List of Supplementary Materials}}

\begin{itemize}[leftmargin=2em,itemsep=2pt]
\item Supplementary Text S1: Ledoit--Wolf shrinkage derivation and
  analytical shrinkage intensity.
\item Supplementary Text S2: Delta method derivation with
  second-order correction terms.
\item Supplementary Text S3: Full Optuna hyperparameter search
  configuration and convergence analysis.
\item Fig.~S1: Reliability diagrams for all six embedding
  representations.
\item Fig.~S2: LOOCV residual distributions by protein family.
\item Fig.~S3: Optuna optimization history (200 trials).
\item Fig.~S4: Embedding t-SNE visualizations colored by $\pKd$.
\item Table~S1: Full LOOCV results for all six representations at
  $N = 24$.
\item Table~S2: Per-pocket breakdown of $\Psuccess$ versus experimental
  activity.
\item Table~S3: Optuna top-10 trial configurations and metrics.
\item Data S1: Processed embeddings and labels for the $N = 932$ dataset
  (NumPy archive).
\item Code S1: BayesDiff source code and pipeline scripts (GitHub
  repository).
\end{itemize}


%% ── Figure captions ──────────────────────────────────────────────────────
\section*{\textbf{Figure Captions}}

\begin{figure}[h!]
\centering
\fbox{\parbox{0.85\textwidth}{\centering\vspace{3cm}
\textbf{Figure 1 placeholder}\\
BayesDiff pipeline overview\vspace{3cm}}}
\caption{\textbf{BayesDiff framework overview.}
A protein pocket is input to the TargetDiff diffusion model, which
generates $M = 64$ molecular conformations.  SE(3)-invariant embeddings
($z \in \mathbb{R}^{128}$) are extracted from the encoder backbone.
Generation uncertainty is estimated via Ledoit--Wolf shrinkage covariance
and GMM multimodal detection.  Prediction uncertainty is obtained from a
Gaussian process with a rational quadratic kernel.  The Delta method
fuses both sources into a total predictive variance $\Utotal$.  Isotonic
regression calibrates the raw probability $\Psuccess$, and Mahalanobis
OOD detection applies a confidence correction, yielding the final
calibrated $\Pfinal \in [0,1]$.}
\label{fig:pipeline}
\end{figure}

\begin{figure}[h!]
\centering
\fbox{\parbox{0.85\textwidth}{\centering\vspace{3cm}
\textbf{Figure 2 placeholder}\\
Representation comparison\vspace{3cm}}}
\caption{\textbf{Representation bottleneck analysis.}
LOOCV Spearman $\rho$ (left) and $R^{2}$ (right) for six molecular
representations on the $N = 932$ dataset.  All 2D fingerprint variants
(ECFP4-128, ECFP4-2048, ECFP6-2048, FCFP4-2048, RDKit-2D) cluster near
$R^{2} \approx 0.01$ (dashed line), whereas the SE(3)-equivariant
Encoder-128 embedding achieves $R^{2} = 0.12$, a $9.2\times$
improvement.  Error bars show 95\% bootstrap confidence intervals.}
\label{fig:representation}
\end{figure}

\begin{figure}[h!]
\centering
\fbox{\parbox{0.85\textwidth}{\centering\vspace{3cm}
\textbf{Figure 3 placeholder}\\
Calibration reliability diagram\vspace{3cm}}}
\caption{\textbf{Calibration reliability diagram.}
Predicted $\Psuccess$ (horizontal axis) versus observed frequency of
activity $y \geq 7.0$ (vertical axis) for the full BayesDiff pipeline
with Encoder-128 on the 48-pocket evaluation set ($\mathrm{ECE} = 0.034$).
The dashed diagonal represents perfect calibration.  Points near the
diagonal indicate well-calibrated probabilities; deviations indicate
over- or under-confidence.}
\label{fig:calibration}
\end{figure}

\begin{figure}[h!]
\centering
\fbox{\parbox{0.85\textwidth}{\centering\vspace{3cm}
\textbf{Figure 4 placeholder}\\
Ablation component contributions\vspace{3cm}}}
\caption{\textbf{Ablation study: component contributions.}
ECE (left) and NLL (right) for the full pipeline and four ablated
variants.  Removing oracle variance (A2) causes the NLL to explode to
${\sim}10^{10}$ (truncated for visibility), demonstrating that $\Uoracle$
is irreplaceable.  Removing generation uncertainty (A1) and
Ledoit--Wolf shrinkage (A4) each produce modest degradations.  The
isotonic calibration ablation (A3) has no effect because insufficient
calibration data are available in the current evaluation setting.}
\label{fig:ablation}
\end{figure}


%% ── Table captions ───────────────────────────────────────────────────────
\section*{\textbf{Table Captions}}

\begin{table}[h!]
\centering
\caption{\textbf{Representation comparison on the initial evaluation set
($N = 24$, LOOCV).}
Six molecular representations evaluated with identical GP models.  All 2D
fingerprint variants produce negative Spearman correlations and negative
or near-zero $R^{2}$, confirming a fundamental representation
bottleneck.  Dim: embedding dimensionality.  95\% CI: coverage of 95\%
predictive intervals.  NaN and --- indicate undefined values due to
divergence or numerical instability.}
\label{tab:rep_n24}
\begin{tabular}{@{}lrrrrc@{}}
\toprule
Embedding & Dim & LOOCV $\rho$ & LOOCV RMSE & LOOCV $R^{2}$ & 95\% CI \\
\midrule
ECFP4-128   &  128  & $-0.37$ & 2.22 & $-0.25$ & 92\% \\
ECFP4-2048  & 2048  & $-0.30$ & 2.43 & ---     & 96\% \\
ECFP6-2048  & 2048  & $-0.32$ & 2.49 & ---     & 96\% \\
FCFP4-2048  & 2048  & $-0.23$ & 2.31 & $-0.36$ & 96\% \\
RDKit-2D    &  217  & $-0.15$ & 3.47 & ---     & 75\% \\
Combined    & 2265  & NaN     & 5.52 & ---     & --- \\
\bottomrule
\end{tabular}
\end{table}

\begin{table}[h!]
\centering
\caption{\textbf{Encoder-128 versus FCFP4-2048 on the expanded dataset
($N = 932$).}
The SE(3)-equivariant encoder embedding achieves a $9.2\times$
improvement in $R^{2}$ and a $3.3\times$ improvement in Spearman $\rho$
over the best 2D fingerprint.  Results from 50$\times$ repeated 70/30
splits confirm robustness.  $p$-values are from two-sided Spearman rank
tests.}
\label{tab:encoder_v_fcfp}
\begin{tabular}{@{}lccc@{}}
\toprule
Metric & Encoder-128 & FCFP4-2048 & Improvement \\
\midrule
LOOCV RMSE      & 1.949  & 2.068  & 5.7\% $\downarrow$ \\
LOOCV $\rho$    & 0.369  & 0.111  & 3.3$\times$ $\uparrow$ \\
LOOCV $R^{2}$   & 0.120  & 0.013  & 9.2$\times$ $\uparrow$ \\
$p$-value       & $<0.0001$ & 0.0007 & --- \\
\midrule
50$\times$ Split $\rho$   & $0.373 \pm 0.045$ & $0.134 \pm 0.055$ & 2.8$\times$ $\uparrow$ \\
50$\times$ Split $R^{2}$  & $0.116 \pm 0.027$ & $0.004 \pm 0.010$ & 29$\times$ $\uparrow$ \\
\bottomrule
\end{tabular}
\end{table}

\begin{table}[h!]
\centering
\caption{\textbf{Ablation study ($N = 24$, Encoder-128, $\ytarget = 7.0$).}
Each row removes one component from the full BayesDiff pipeline.
Removing oracle variance (A2) causes the NLL to explode, demonstrating
that $\Uoracle$ is essential.  ECE: expected calibration error; AUROC:
area under the ROC curve; NLL: negative log-likelihood.}
\label{tab:ablation}
\begin{tabular}{@{}llccc@{}}
\toprule
Variant & Description & ECE & AUROC & NLL \\
\midrule
Full    & Complete pipeline       & 0.264 & 1.000 & 1.95 \\
A1      & No $\Ugen$ ($\Siggen = 0$) & 0.230 & 1.000 & 1.95 \\
A2      & No $\Uoracle$ ($\Uoracle = 0$) & 0.238 & 1.000 & $1.03 \times 10^{10}$ \\
A3      & No calibration          & 0.264 & 1.000 & 1.95 \\
A4      & Naive covariance (no L--W) & 0.275 & 1.000 & 1.98 \\
\bottomrule
\end{tabular}
\end{table}


\end{document}
